\begin{figure}[H]
\centering
\begin{tikzpicture}[>=Stealth,scale=0.91]
  \draw[rounded corners,very thick,gray!65] (-5,-2.35) rectangle (5.5,2.35);
  \draw[->] (-5.4,0) -- (7.2,0) node[right] {eje óptico};
  \draw[very thick,blue!65] (-4.35,-2.05) .. controls (-4.1,-1.2) and (-4.6,-0.5) .. (-4.35,0)
    .. controls (-4.1,0.5) and (-4.6,1.2) .. (-4.35,2.05);
  \node[rotate=90] at (-4.82,0) {placa correctora};
  \draw[very thick] (4.5,-2.0) .. controls (5.15,-1) and (5.15,-0.35) .. (4.55,-0.28);
  \draw[very thick] (4.55,0.28) .. controls (5.15,0.35) and (5.15,1) .. (4.5,2.0);
  \node[right] at (4.85,-1.55) {primario esférico};
  \draw[very thick] (-4.05,-0.58) .. controls (-3.55,-0.3) and (-3.55,0.3) .. (-4.05,0.58);
  \node at (-2.0,1.85) {secundario convexo};
  \draw[->,thin] (-2.45,1.68) -- (-3.8,0.42);
  \foreach \y in {-1.3,1.3}{
    \draw[->,blue,thick] (-5.0,\y) -- (4.62,\y);
  }
  \draw[red,thick] (4.62,1.3) -- (-3.8,0.4);
  \draw[red,thick] (4.62,-1.3) -- (-3.8,-0.4);
  \draw[->,green!60!black,thick] (-3.8,0.4) -- (6.2,0);
  \draw[->,green!60!black,thick] (-3.8,-0.4) -- (6.2,0);
  \filldraw[fill=gray!25] (6.1,-0.5) rectangle (6.35,0.5);
  \node[above] at (6.22,0.58) {plano focal};
\end{tikzpicture}
\caption{Elementos ópticos de un telescopio Schmidt--Cassegrain.}
\end{figure}